\section{Social-Science Measurement Setting}\label{sec:social-measurement}

The semantic application in the next section is not just a long-context
summarization benchmark. It is a measurement problem: a researcher has a
rubric, a trusted but expensive target judgment, and a downstream analysis
that can be biased if cheaper labels preserve surface accuracy while
dropping task-relevant distinctions. That is the setting in which the
oracle language of Sections~\ref{sec:framework}--\ref{sec:theorems}
matters most.

\subsection{Rubrics, Expert Means, and Agreement}

Comparative-politics text measurement is built around expert-defined
constructs: economic left--right placement, social liberalism, immigration,
European integration, environmental policy, decentralization, and related
issue scales. The Manifesto Project provides the long-document corpus
\citep{ManifestoProject2025a}; \citet{BenoitEtAl2025} provide the recent
LLM benchmark we use, comparing model-produced scores to expert-survey
means on seven-point policy scales. In this design, the target is not a
single annotator's sentence label. It is a construct-level party-position
estimate obtained by aggregating expert judgments, and split-expert
correlations are an empirical agreement reference for how noisy that
measurement target is.

This matters for C-TreePO because the object being preserved is the
rubric-indexed judgment, not the literal text. A summary can safely discard
stylistic details, slogans, and redundant examples if the expert-scale
position remains unchanged. It fails when it drops evidence that changes
the issue placement or a cross-section interaction that an expert would
have used to resolve an ambiguous platform.

\subsection{Surrogate Labels and Downstream Inference}

LLMs are attractive in this setting because they can cheaply produce
rubric labels, explanations, and summaries for many long documents.
But social-science use is not finished when a prediction correlates with
expert means. \citet{EgamiEtAl2024} show that using imperfect surrogate
labels in downstream inference can induce bias even when prediction
quality is high, and their design-based supervised-learning framework
uses gold-standard annotations to correct that risk. C-TreePO addresses
the complementary compression question: when a long document is replaced
by a tree summary, what evidence has to survive for the downstream
oracle-indexed claim to remain valid?

The audit in Section~\ref{sec:estimation} is therefore a measurement
instrument, not a generic quality score. It samples leaf, idempotence, and
merge checks from the realized tree; the calibration residual measures the
gap between the judge used during training or auditing and the expert
oracle the researcher ultimately trusts; and the finite-sample envelope
states how much uncertainty remains after collecting global and local
labels.

\subsection{Why the Manifesto Benchmark Fits}

Manifestos have exactly the structure the framework is meant to test.
They are long, sectioned documents; their issue positions are rubric-coded;
their evidence is often local but can be resolved by interactions across
sections; and the benchmark reports external expert means rather than only
model-internal consistency. Section~\ref{sec:manifesto-llm} keeps the
empirical protocol and numbers in one place. The interpretation, fixed
here, is that matching the benchmark is evidence that a single
open-weight, tree-structured prompt program can preserve a
social-science measurement target, while the local-law audit supplies
the failure modes and finite-sample caveats that a flat long-context
read does not expose.
